\documentclass[../../main.tex]{subfiles}
\begin{document}

\begin{center}
  {\fontsize{48pt}{58pt}\selectfont \textbf{1982}\normalsize\textbf{年度}}
\end{center}

\begin{table}[htb]
  \centering
  \caption{出題テーマと難易度}
  \begin{tabular}{|c|p{6.2cm}|c|c|c|}
    \hline
    \textbf{問題番号} & \textbf{テーマ} & \textbf{思考量} & \textbf{計算量} & \textbf{総合} \\
    \hline
    第1問 & 一次変換の不動点と不変直線 &  &  &  \\
    \hline
    第2問 & 正四面体と辺に接する球 & 普 & 普 & 普 \\
    \hline
    第3問 & 円を横切る線分の長さ & 普 & 普 & 普 \\
    \hline
    第4問 & 正弦曲線上の等速運動と加速度 & 普 & 難 & 普 \\
    \hline
    第5問 & 不等式で表される立体の体積 & 易 & 易 & 易 \\
    \hline
    第6問 & サイコロの回転と状態遷移 & 普 & 普 & 普 \\
    \hline
  \end{tabular}
\end{table}

\begin{table}[htb]
  \centering
  \caption{解答}
  \renewcommand{\arraystretch}{1.4}
  \begin{tabular}{|c|c|p{9cm}|}
    \hline
    \textbf{問題番号} & \textbf{小問} & \textbf{答え} \\
    \hline
    第1問 & -- & \\
    \hline
    第2問 & -- & 一辺の長さは$2\sqrt{2}$．体積は$\dfrac{8\pi}{3}\left(1-\dfrac{4\sqrt{3}}{9}\right)$ \\
    \hline
    第3問 & (1) & $\dfrac{\pi}{4}<\theta<\dfrac{\pi}{3}$ \\
    \cline{2-3}
    & (2) & $BC=3-4\sin^2\theta$ \\
    \cline{2-3}
    & (3) & $\dfrac{\sqrt{14}}{8}\le CM<\dfrac{\sqrt{2}}{2}$ \\
    \hline
    第4問 & -- & 加速度の大きさの最大値は$V^2$ \\
    \hline
    第5問 & -- & $V=\dfrac{7}{4}$ \\
    \hline
    第6問 & (1) & $(p_1,q_1,r_1)=\left(\dfrac{1}{3},\dfrac{2}{3},0\right)$ \\
    \cline{2-3}
    & (2) & $p_n=\dfrac{p_{n-1}}{3}+\dfrac{q_{n-1}}{6}$，$q_n=\dfrac{2}{3}$，$r_n=\dfrac{q_{n-1}}{6}+\dfrac{r_{n-1}}{3}$ \\
    \cline{2-3}
    & (3) & $(p,q,r)=\left(\dfrac{1}{6},\dfrac{2}{3},\dfrac{1}{6}\right)$ \\
    \hline
  \end{tabular}
\end{table}

\paragraph{試験形態}
6問．試験時間は要確認．

\end{document}
